% DEFINITIONS: Formal Conjectures for Deep Number-Theoretic Foundations
% of Ultrametric Quantum Information
% Project: number-theory-ultrametric-deep
% Date: 2026-07-03
% Status: Phase 1 — Conjecture Formalization
\documentclass[12pt]{article}
\usepackage{amsmath, amssymb, amsthm, mathtools}
\usepackage{mathrsfs}
\usepackage[utf8]{inputenc}
\usepackage[T1]{fontenc}

\newtheorem{definition}{Definition}[section]
\newtheorem{conjecture}[definition]{Conjecture}
\newtheorem{question}[definition]{Open Question}
\theoremstyle{remark}
\newtheorem{remark}[definition]{Remark}
\newtheorem{notation}[definition]{Notation}

\title{Formal Definitions and Conjectures for Deep Number-Theoretic \\
Quantum Error Correction}
\author{QNFO Research Collective}
\date{2026-07-03}

\begin{document}
\maketitle

\begin{abstract}
This document formalizes the central definitions and conjectures connecting
seven pillars of modern number theory to quantum error correction, hierarchical
discovery, and ultrametric memory architectures. Each definition specifies a
proposed mathematical object or structure. Each conjecture states a relationship
to be proved or falsified. Certainty labels follow QNFO-POL-COM-001:
$[\text{established}]$, $[\text{conjecture}]$, $[\text{speculative}]$,
$[\text{not yet falsifiable}]$. The document serves as the formal reference
for Phase 2 computational exploration and Phase 3 theorem development.
\end{abstract}

% ======================================================================
\section{Preliminaries and Notation}
% ======================================================================

\begin{notation}[Stabilizer Codes]
A \emph{stabilizer code} $C = [[n,k,d]]_q$ is a $q^k$-dimensional subspace of
$(\mathbb{C}^q)^{\otimes n}$ defined as the simultaneous $+1$ eigenspace of an
abelian subgroup $\mathcal{S} \subset \mathcal{P}_n(q)$, the $n$-qudit Pauli
group, where $-\mathbb{I} \notin \mathcal{S}$ and $|\mathcal{S}| = q^{n-k}$.
The code distance $d = \min_{P \in N(\mathcal{S}) \setminus \mathcal{S}}
\operatorname{wt}(P)$ where $N(\mathcal{S})$ is the normalizer of $\mathcal{S}$
in $\mathcal{P}_n(q)$.
\end{notation}

\begin{notation}[p-adic Numbers]
For a prime $p$, let $\mathbb{Q}_p$ denote the field of $p$-adic numbers with
valuation $v_p$ and absolute value $|x|_p = p^{-v_p(x)}$. Let $\mathbb{Z}_p$
denote the ring of $p$-adic integers. Let $\mathbb{C}_p$ denote the completion
of an algebraic closure of $\mathbb{Q}_p$. The valuation extends uniquely to
$\mathbb{C}_p$.
\end{notation}

\begin{notation}[Witt Vectors]
Let $W(\mathbb{F}_q)$ denote the ring of Witt vectors of the finite field
$\mathbb{F}_q$ ($q = p^m$). Elements are infinite sequences
$(a_0, a_1, a_2, \ldots)$ with $a_i \in \mathbb{F}_q$, with addition and
multiplication defined by the universal Witt polynomials.
$W_n(\mathbb{F}_q)$ denotes the truncated Witt vectors of length $n$.
\end{notation}

\begin{notation}[Galois Groups]
$G_{\mathbb{Q}_p} = \operatorname{Gal}(\overline{\mathbb{Q}}_p / \mathbb{Q}_p)$
is the absolute Galois group of $\mathbb{Q}_p$. $I_p \subset G_{\mathbb{Q}_p}$
is the inertia subgroup. $W_{\mathbb{Q}_p}$ is the Weil group.
\end{notation}

% ======================================================================
\section{Pillar I: p-adic Analysis of Quantum Codes}
% ======================================================================

\begin{definition}[p-adic Encoding of a Stabilizer Code]
$[\text{definition}]$
Let $C = [[n,k,d]]_q$ be a stabilizer code over $\mathbb{C}$. A
\emph{p-adic encoding} of $C$ at prime $p$ is a continuous map
$\iota_C: \mathbb{Z}_p \to \mathbb{Q}_p$ constructed from the code
parameters $(n,k,d)$ and a choice of basis such that:
\begin{enumerate}
\item $\iota_C$ is continuous (i.e., uniformly continuous w.r.t.\ the
  $p$-adic metric)
\item The Mahler expansion $\iota_C(x) = \sum_{m=0}^{\infty} a_m \binom{x}{m}$
  converges uniformly, with $a_m \in \mathbb{Q}_p$ and $|a_m|_p \to 0$
\item The coefficient sequence $(a_m)_{m \geq 0}$ is computable from $C$
  in polynomial time in $n$
\end{enumerate}
The encoding is \emph{natural} if it is invariant under local Clifford
transformations of $C$ up to a $p$-adic isometry.
\end{definition}

\begin{conjecture}[Mahler Decay Classifies Code Structure]
$[\text{conjecture}]$
The Mahler coefficient decay rate $\rho(C) = -\limsup_{m \to \infty}
\frac{\log |a_m|_p}{\log m}$ is an invariant of the stabilizer code $C$
(independent of the choice of natural encoding) and satisfies:
\[
\rho(C) \begin{cases}
= 0 & \text{if } C \text{ is a random CSS code} \\
\in (0, 1] & \text{if } C \text{ is an algebraic code (Goppa, AG)} \\
\text{correlated with } d/n & \text{for surface codes}
\end{cases}
\]
This would be disconfirmed if two natural encodings of the same code produce
statistically distinguishable Mahler spectra.
\end{conjecture}

\begin{question}[Amice Spectral Classification]
$[\text{speculative}]$
Does the intrinsic Amice transform of $\iota_C$ detect Galois-invariant
subspaces in the space of quantum codes? Specifically, for Goppa codes
(which arise from curves over finite fields), is there a characteristic
spectral signature in the Amice domain that distinguishes them from
random CSS codes of the same $[[n,k,d]]$ parameters?
\end{question}

% ======================================================================
\section{Pillar II: Hasse Local-Global Principle}
% ======================================================================

\begin{definition}[Local Quantum Code at Prime $p$]
$[\text{definition — central to entire framework}]$
Let $C = [[n,k,d]]_q$ be a stabilizer code. For a prime $p$, a
\emph{local quantum code at $p$} with the same parameters $[[n,k,d]]_q$
is a stabilizer code $C_p$ defined over the $p$-adic numbers satisfying:
\begin{enumerate}
\item $C_p$ is defined by a stabilizer group $\mathcal{S}_p$ whose matrix
  entries (in some $p$-adic basis) lie in $\mathbb{Q}_p$ (or its
  unramified extension)
\item The code subspace is a $q^k$-dimensional subspace of a $p$-adic
  Hilbert module $(\mathbb{Q}_p(\zeta_q))^{\otimes n}$ where
  $\zeta_q = e^{2\pi i/q}$
\item The distance $d_p$ (minimum $p$-adic weight of non-trivial logical
  operators) satisfies $|d_p - d|_p < 1$ (i.e., $d$ is the reduction of
  $d_p$ modulo $p$)
\item There exists a \emph{Witt lift} from characteristic $p$: a stabilizer
  code $\overline{C}_p$ over $\mathbb{F}_p$ (or $\mathbb{F}_{p^m}$) whose
  Witt vector lift $W(\overline{C}_p) \otimes_{\mathbb{Z}_p} \mathbb{Q}_p$
  is locally equivalent to $C_p$
\end{enumerate}
\end{definition}

\begin{remark}
This definition is the SINGLE most important formalization in the
entire research plan. Without a rigorous notion of ``local quantum code''
at a prime, Pillars II, III, VI, and VII have no objects to work with.
The Witt lift condition (item 4) connects this to Pillar VI (Dieudonn\'e
modules) and Pillar III (crystalline codes).
\end{remark}

\begin{conjecture}[Hasse Principle for Stabilizer Codes]
$[\text{conjecture}]$
A stabilizer code with parameters $[[n,k,d]]_q$ exists over $\mathbb{C}$
if and only if it exists locally at EVERY prime $p$ AND at the archimedean
place (i.e., over $\mathbb{R}$ or $\mathbb{C}$). Formally:
\[
\exists C/\mathbb{C} \iff
\left(\forall p \text{ prime}: \exists C_p/\mathbb{Q}_p\right) \land
\left(\exists C_\infty/\mathbb{R}\right)
\]
This would be disconfirmed if code parameters are found that are everywhere
locally realizable but not globally realizable (analogous to Selmer's
counterexample $3x^3 + 4y^3 + 5z^3 = 0$ for the classical Hasse principle).
\end{conjecture}

\begin{conjecture}[Brauer-Manin Obstruction for Quantum Codes]
$[\text{conjecture}]$
When the Hasse principle fails, the obstruction is captured by a
cohomological invariant taking values in a Brauer-like group
$\operatorname{Br}(\mathcal{C})$ associated to the ``moduli space''
of stabilizer codes, analogous to $\operatorname{Br}(X)/\operatorname{Br}(k)$
for algebraic varieties.
\end{conjecture}

\begin{question}[Minkowski-Hasse for Codes]
$[\text{speculative}]$
Does there exist a classification theorem for stabilizer code parameters
analogous to the Minkowski-Hasse theorem for quadratic forms — i.e., a
complete set of local invariants (dimension, $p$-adic signature, Witt index)
that determine when parameters $[[n,k,d]]_q$ are realizable?
\end{question}

% ======================================================================
\section{Pillar III: Galois Representations and p-adic Hodge Theory}
% ======================================================================

\begin{definition}[Galois Representation of a Local Quantum Code]
$[\text{definition — speculative}]$
Let $C_p$ be a local quantum code at prime $p$ (as defined in \S 2).
A \emph{Galois representation} associated to $C_p$ is a continuous
homomorphism
\[
\rho_{C_p}: G_{\mathbb{Q}_p} \to \operatorname{GL}_n(\mathbb{Q}_p)
\]
constructed from the stabilizer group $\mathcal{S}_p$ and the Witt lift
$\overline{C}_p$, satisfying:
\begin{enumerate}
\item $\rho_{C_p}$ is unramified if and only if $C_p$ lifts from an
  unramified extension of $\mathbb{Q}_p$ (crystalline case)
\item $\det \circ \rho_{C_p}$ is the $p$-adic cyclotomic character times
  the code rate $k/n$
\item The representation is Hodge-Tate with weights corresponding to the
  code distance spectrum
\end{enumerate}
\end{definition}

\begin{conjecture}[Fontaine Filtration of Stabilizer Codes]
$[\text{conjecture}]$
There exists a filtration of the category of local quantum codes at $p$:
\[
\mathcal{C}_{\text{crys}} \subset
\mathcal{C}_{\text{st}} \subset
\mathcal{C}_{\text{dR}}
\]
corresponding to Fontaine's period ring filtration
$B_{\text{crys}} \subset B_{\text{st}} \subset B_{\text{dR}}$, where:
\begin{itemize}
\item $\mathcal{C}_{\text{crys}}$ (crystalline codes): lift from $\mathbb{F}_p$
  without ramification; codes with ``good reduction'' — distance preserved
  under $p$-adic deformation
\item $\mathcal{C}_{\text{st}}$ (semistable codes): allow controlled
  ramification; Clifford-deformed stabilizers
\item $\mathcal{C}_{\text{dR}}$ (de Rham codes): all codes expressible
  in the $p$-adic analytic category
\end{itemize}
\end{conjecture}

\begin{conjecture}[Fontaine-Mazur Analog for Quantum Codes]
$[\text{speculative}]$
A quantum code $C$ is ``geometric'' (physically realizable by physical qubits
in a laboratory) if and only if its associated Galois representation
$\rho_{C_p}$ is de Rham (in the sense of Fontaine) and unramified at almost
all primes. This would provide a mathematical criterion for physical
realizability independent of hardware constraints.
\end{conjecture}

\begin{question}[Frobenius = Cyclic Measurement]
$[\text{speculative}]$
The Frobenius endomorphism $\varphi$ on the $p$-adic Galois representation
$\rho_{C_p}$ satisfies $\varphi(x) \equiv x^p \pmod{p}$. In the silent-radix
framework, cyclic measurement operators generate finite cyclic groups on
quantum state spaces. Does the Frobenius action on $\rho_{C_p}$ correspond
to a cyclic measurement operator? If so, what is the operational meaning
of the characteristic polynomial of Frobenius?
\end{question}

% ======================================================================
\section{Pillar IV: Class Field Theory and Cyclic Measurement}
% ======================================================================

\begin{definition}[Artin Symbol of a Cyclic Measurement]
$[\text{definition}]$
Let $\mathcal{M}_D$ be a cyclic measurement protocol with period $D$
(silent-radix framework, $D=4$ Fourier clocks). The group generated by
$\mathcal{M}_D$ is isomorphic to $\mathbb{Z}/D\mathbb{Z}$. The
\emph{Artin symbol} $\operatorname{Art}(\mathcal{M}_D)$ is the element
of $\operatorname{Gal}(\mathbb{Q}_p(\zeta_D)/\mathbb{Q}_p)$
corresponding to the measurement operator under the local Artin map:
\[
\theta_p: \mathbb{Q}_p^\times \to
\operatorname{Gal}(\mathbb{Q}_p^{\text{ab}}/\mathbb{Q}_p)
\]
via the identification $\mathbb{Z}/D\mathbb{Z} \cong
(\mathbb{Z}/D\mathbb{Z})^\times$ (for $D$ coprime to $p$).
\end{definition}

\begin{conjecture}[Kronecker-Weber Theorem for Quantum Measurement]
$[\text{conjecture}]$
Every abelian measurement scheme (every finite abelian group of unitary
operators acting on quantum state space that forms a complete set of
measurements) is contained in a generalized Clifford group $\mathcal{C}_n(q)$,
just as every finite abelian extension of $\mathbb{Q}$ is contained in a
cyclotomic extension $\mathbb{Q}(\zeta_m)$ (Kronecker-Weber theorem).
This would be disconfirmed by exhibiting an abelian measurement scheme
not contained in any generalized Clifford group.
\end{conjecture}

\begin{question}[Idele Class Group = Universal Measurement Group]
$[\text{speculative}]$
Does the idele class group $\mathbb{A}_\mathbb{Q}^\times/\mathbb{Q}^\times$
serve as a universal measurement group for the adelic QEC framework,
where each local component $\mathbb{Q}_p^\times$ corresponds to measurements
at the $p$-adic completion?
\end{question}

% ======================================================================
\section{Pillar V: Arithmetic Geometry of Quantum Codes}
% ======================================================================

\begin{definition}[p-adic Deformation Family of Quantum Codes]
$[\text{definition}]$
A \emph{$p$-adic deformation family} of stabilizer codes is a 1-parameter
family $\{C_t\}_{t \in \mathbb{Z}_p}$ such that:
\begin{enumerate}
\item For each $t \in \mathbb{Z}_p$, $C_t$ is a stabilizer code with
  parameters $[[n_t, k_t, d_t]]_q$
\item The stabilizer generators vary continuously (in the $p$-adic topology)
  as functions of $t$
\item The special fiber $C_0$ (at $t=0$) is the ``reference code''
\item The distance function $t \mapsto d_t$ is locally constant on $p$-adic
  discs (ultrametric property)
\end{enumerate}
\end{definition}

\begin{conjecture}[Kodaira-N\'eron Classification of Code Degenerations]
$[\text{conjecture}]$
Let $\{C_t\}_{t \in \mathbb{Z}_p}$ be a $p$-adic deformation family of
quantum codes. The behavior of the distance $d_t$ as $t \to 0$ (i.e., under
$p$-adic reduction) follows the Kodaira-N\'eron classification:
\begin{itemize}
\item \textbf{Good reduction} (Type I$_0$): $d_t = d_0$ for all $t$ with
  $|t|_p$ sufficiently small — distance is preserved
\item \textbf{Multiplicative reduction} (Type I$_n$): $d_t \to d_0/n$ as
  $t \to 0$ — partial distance collapse
\item \textbf{Additive reduction} (Types II-IV$^*$): $d_t \to 0$ —
  catastrophic distance collapse
\end{itemize}
\end{conjecture}

\begin{question}[Isogeny Graph = Code-Switching Protocol]
$[\text{speculative}]$
Does a path in a supersingular isogeny graph between two elliptic curves
$E_1$ and $E_2$ correspond to a sequence of code deformations (a
``code-switching protocol'') that preserves ultrametric invariants while
transforming code parameters? If so, the Ramanujan expander property of
isogeny graphs would guarantee efficient code-switching.
\end{question}

% ======================================================================
\section{Pillar VI: Witt Vectors and Stabilizer Code Cohomology}
% ======================================================================

\begin{definition}[Dieudonn\'e Module of a Stabilizer Code]
$[\text{definition — speculative}]$
Let $\overline{C}$ be a stabilizer code over $\mathbb{F}_p$ (the reduction
mod $p$ of a local quantum code $C_p$). A \emph{Dieudonn\'e module}
associated to $\overline{C}$ is a module $D(\overline{C})$ over the
Dieudonn\'e ring $D_{\mathbb{F}_p} = W(\mathbb{F}_p)[F, V]/(FV - p)$
where:
\begin{itemize}
\item $F: D(\overline{C}) \to D(\overline{C})$ is a semilinear operator
  (Frobenius) satisfying $F(ax) = \sigma(a) F(x)$ where $\sigma$ is the
  Frobenius automorphism of $W(\mathbb{F}_p)$
\item $V: D(\overline{C}) \to D(\overline{C})$ is a semilinear operator
  (Verschiebung) satisfying $FV = VF = p$
\end{itemize}
The module $D(\overline{C})$ encodes the cohomological structure of the
stabilizer code, with the first cohomology group $H^1_{\text{stab}}(C)$
being isomorphic to $D(\overline{C})$ as a $D_{\mathbb{F}_p}$-module.
\end{definition}

\begin{conjecture}[Frobenius = Clifford Conjugation, Verschiebung = T-gate]
$[\text{conjecture}]$
For a stabilizer code $\overline{C}$ over $\mathbb{F}_p$:
\begin{enumerate}
\item The Frobenius operator $F$ on $D(\overline{C})$ corresponds to
  conjugation by an element of the Clifford group $\mathcal{C}_n(p)$
\item The Verschiebung operator $V$ corresponds to T-gate injection
  (a non-Clifford operation that increases the level of the Clifford
  hierarchy)
\item The relation $FV = p = VF$ corresponds to the fact that $p$
  iterations of the T-gate compose to a Clifford operation (for qudits
  of dimension $p$)
\end{enumerate}
This would be disconfirmed if no natural $F$-$V$ pair can be constructed
on the stabilizer code space satisfying the Dieudonn\'e relation and
respecting the quantum computational interpretation.
\end{conjecture}

\begin{conjecture}[Dieudonn\'e-Manin Slope = Error-Correction Capacity]
$[\text{conjecture}]$
The Dieudonn\'e-Manin classification decomposes the isocrystal
$D(\overline{C}) \otimes_{\mathbb{Z}_p} \mathbb{Q}_p$ into isoclinic
components by their slopes $\lambda \in \mathbb{Q}$. Then:
\begin{itemize}
\item \textbf{Crystalline codes} (integer slopes): sharp error-correction
  thresholds — the code distance is a rigid invariant
\item \textbf{Non-crystalline codes} (fractional slopes): soft thresholds,
  distance degrades continuously with noise
\item The slope $\lambda$ is monotonically related to the threshold error
  rate: higher slopes $\implies$ lower thresholds
\end{itemize}
\end{conjecture}

\begin{conjecture}[Witt Truncation = Thermodynamic Limit]
$[\text{conjecture}]$
Truncated Witt vectors $W_n(\mathbb{F}_q)$ correspond to
$\mathbb{Z}/p^n\mathbb{Z}$-truncated quantum codes (codes with $n$ levels
of concatenation). The inverse limit $W(\mathbb{F}_q) =
\varprojlim W_n(\mathbb{F}_q)$ corresponds to the thermodynamic limit
$n \to \infty$ of infinite code families, with the Witt vector
$p$-adic topology encoding the convergence of code families.
\end{conjecture}

% ======================================================================
\section{Pillar VII: Bruhat-Tits Buildings, Berkovich Spaces, Langlands}
% ======================================================================

\begin{definition}[Moy-Prasad Filtration of Quantum Codes]
$[\text{definition — speculative}]$
Let $C$ be a quantum code embedded (via p-adic encoding) into a smooth
representation $\pi_C$ of $\operatorname{GL}_n(\mathbb{Q}_p)$ or an
associated $p$-adic group. The \emph{Moy-Prasad depth} $r(C)$ is the
minimal $r \geq 0$ such that $\pi_C$ has nonzero $G(\mathbb{Q}_p)_{x,r}$-
fixed vectors for some point $x$ in the Bruhat-Tits building
$\mathcal{B}(\operatorname{GL}_n, \mathbb{Q}_p)$.
\end{definition}

\begin{conjecture}[Moy-Prasad Depth = Logical Error Rate]
$[\text{conjecture}]$
The Moy-Prasad depth $r(C)$ is related to the logical error rate
$p_L(C, \eta)$ of the code $C$ under a noise model with physical error
rate $\eta$:
\[
p_L(C, \eta) \sim \exp(-\alpha \cdot r(C) \cdot d)
\]
for some constant $\alpha > 0$. Codes with larger depth have lower error
rates; the depth corresponds to the fault-tolerance threshold in the sense
that codes with $r(C) > r_{\text{crit}}$ are below threshold.
\end{conjecture}

\begin{conjecture}[Bernstein Decomposition = Universality Classes]
$[\text{conjecture}]$
The Bernstein decomposition of the category of smooth representations of
$G(\mathbb{Q}_p)$ into Bernstein blocks classifies universality classes
of quantum codes: codes in the same Bernstein block share the same
asymptotic scaling behavior (threshold, overhead, decoding complexity).
Different blocks correspond to different code families (surface codes,
hyperbolic codes, LDPC codes, etc.).
\end{conjecture}

\begin{conjecture}[Local Langlands for Quantum Codes]
$[\text{speculative — deferred}]$
The local Langlands correspondence for $\operatorname{GL}_n(\mathbb{Q}_p)$:
\[
\begin{Bmatrix}
\text{irreducible smooth representations} \\
\text{of } \operatorname{GL}_n(\mathbb{Q}_p)
\end{Bmatrix}
\longleftrightarrow
\begin{Bmatrix}
n\text{-dimensional Weil-Deligne} \\
\text{representations of } W_{\mathbb{Q}_p}
\end{Bmatrix}
\]
classifies $n$-qudit quantum codes: the Langlands parameter
$(\rho, N)$ encodes the code's weight enumerator
($L$-function analog) and error statistics
($\varepsilon$-factor analog). The qudit count $n$ corresponds to
$\operatorname{GL}_n$, and supercuspidal representations correspond
to primitive (indecomposable) quantum codes.
\end{conjecture}

\begin{remark}
This is the most ambitious conjecture in the framework. Even the
classical local Langlands correspondence for $\operatorname{GL}_n$
is a deep theorem (Harris-Taylor 2001, Henniart 2000). Extending it
to quantum codes is a multi-decade research program. The conjecture
is labeled $[\text{speculative — deferred}]$ and should not be a
Phase 2-3 deliverable.
\end{remark}

% ======================================================================
\section{Cross-Cutting Conjectures}
% ======================================================================

\begin{conjecture}[Adelic Code Product Formula]
$[\text{conjecture}]$
Let $C$ be a global quantum code (over $\mathbb{C}$) with local realizations
$C_p$ at each prime $p$ and $C_\infty$ at the archimedean place. Then the
product of local code distances satisfies:
\[
\prod_{p \leq \infty} d_p(C) = 1
\]
(under an appropriate normalization of distance). This would be the quantum
code analog of the Artin product formula $\prod_v |x|_v = 1$ for
$x \in \mathbb{Q}^\times$, imposing a fundamental constraint on how code
distance can vary across completions.
\end{conjecture}

\begin{conjecture}[p-adic to Archimedean Bridge via Fontaine-Fargues]
$[\text{not yet falsifiable}]$
The Fontaine-Fargues curve — a ``$p$-adic Riemann surface'' interpolating
between characteristic 0 and characteristic $p$ — provides a geometric
framework for connecting $p$-adic quantum codes (characteristic $p$ side)
to physical quantum codes (characteristic 0 / archimedean side). The
archimedean limit of ultrametric quantum codes corresponds to the
semiclassical limit $p \to \infty$.
\end{conjecture}

\begin{question}[Arithmetic Topology and Topological Codes]
$[\text{speculative}]$
Arithmetic topology (Mazur, Kapranov, Reznikov) draws analogies:
\[
\{\text{knots in } S^3\} \longleftrightarrow
\{\text{primes in } \operatorname{Spec}(\mathbb{Z})\}
\]
If topological quantum codes (surface codes) are defined on 3-manifolds
via knot theory, does arithmetic topology suggest a notion of
\emph{arithmetic quantum codes} defined on $\operatorname{Spec}(\mathcal{O}_K)$
for a number field $K$, where the genus of $K$ controls the code rate
and prime ideals correspond to excitations?
\end{question}

% ======================================================================
\section{Falsifiability Conditions}
% ======================================================================

\begin{center}
\begin{tabular}{|l|p{8cm}|}
\hline
\textbf{Pillar} & \textbf{Falsification Condition} \\
\hline
I (p-adic analysis) & No correlation between Mahler coefficient decay
  and code structure exists in real benchmark data \\
\hline
II (Hasse principle) & All known code parameters that exist locally also
  exist globally (principle is tautological) \\
\hline
III (Galois reps) & Crystalline/semistable/de Rham provides no
  discrimination between known code types \\
\hline
IV (Class field theory) & Cyclic measurement operators do not form a
  group isomorphic to a Galois group \\
\hline
V (Arithmetic geometry) & Kodaira-N\'eron types show no correlation with
  code property variation under deformation \\
\hline
VI (Witt cohomology) & No natural $F$-$V$ pair exists on the space of
  stabilizer codes \\
\hline
VII (BT/Langlands) & No meaningful map between Moy-Prasad depth and
  code parameters exists \\
\hline
\end{tabular}
\end{center}

% ======================================================================
\section{Minimal Pilot Experiments (Phase 2 Entry Points)}
% ======================================================================

\begin{enumerate}
\item \textbf{Mahler spectrum of 5-qubit code} (Pillar I):
  Compute Mahler coefficients for a specific p-adic encoding of the
  5-qubit perfect code. Compare against random p-adic functions.
  \textit{Cost: 1 day. Success: statistically distinguishable spectra.}

\item \textbf{Local invariants for CSS codes} (Pillar II):
  Define and compute local invariants for small CSS code families
  (toric codes, Shor codes) across small primes $p = 2, 3, 5$.
  \textit{Cost: 2 days. Success: at least one code family shows
  nontrivial variation across primes.}

\item \textbf{Artin symbol for $D=4$ Fourier clocks} (Pillar IV):
  Map the cyclic measurement group $\mathbb{Z}/4\mathbb{Z}$ from
  silent-radix to a specific element of
  $\operatorname{Gal}(\mathbb{Q}_p(\zeta_4)/\mathbb{Q}_p)$.
  \textit{Cost: 1 day. Success: identifiable group isomorphism.}
\end{enumerate}

\end{document}
